% Phase 2: 对称性分类实验结果 (LaTeX)
% 可以直接插入到论文中

\section{对称性分类实验}

\subsection{实验设置}

本实验旨在训练一个独立的对称性分类网络，用于判断3D点云模型的对称性类型。

\subsubsection{数据集}
\begin{itemize}
    \item \textbf{数据集}: ModelNet40 (对齐版本)
    \item \textbf{总标注数}: 378个样本
    \item \textbf{输入}: 10000点点云 + 竖直方向向量 $\mathbf{v}_{\text{up}} = [0, 1, 0]^T$
    \item \textbf{输出}: 5分类 ($K \in \{-1, 0, 1, 2, 4\}$)
\end{itemize}

\subsubsection{数据划分}
采用固定随机种子 (seed=42) 进行分层采样：
\begin{itemize}
    \item \textbf{训练集}: 262个样本 (69.3\%) - 使用旋转增强
    \item \textbf{验证集}: 55个样本 (14.6\%) - 不使用增强
    \item \textbf{测试集}: 61个样本 (16.1\%) - 不使用增强
\end{itemize}

\subsection{标注数据分布}

\begin{table}[h]
\centering
\caption{对称性类型分布}
\label{tab:symmetry_distribution}
\begin{tabular}{clrrr}
\hline
\textbf{K值} & \textbf{含义} & \textbf{训练集} & \textbf{验证集} & \textbf{测试集} \\
\hline
$K=-1$ & 没有正面 & 0 & 0 & 0 \\
$K=0$ & 完全对称 & 1 & 0 & 1 \\
$K=1$ & 1个正面 & 72 & 15 & 17 \\
$K=2$ & 2个正面 & 0 & 0 & 1 \\
$K=4$ & 4个正面 & 189 & 40 & 42 \\
\hline
\textbf{总计} & & 262 & 55 & 61 \\
\hline
\end{tabular}
\end{table}

\subsection{消融实验设计}

\subsubsection{实验1: Upright向量的输入方式}

我们设计了三种网络架构来验证upright向量的最佳输入方式：

\begin{itemize}
    \item \textbf{Exp 1.1 (GlobalConcat)}:
        将upright向量在全局特征后拼接
        $$\mathbf{f}_{\text{final}} = [\mathbf{f}_{\text{global}}; \mathbf{v}_{\text{up}}] \in \mathbb{R}^{1027}$$

    \item \textbf{Exp 1.2 (PointConcat)}:
        将upright向量复制到每个点作为额外特征
        $$\mathbf{p}_i' = [\mathbf{p}_i; \mathbf{v}_{\text{up}}] \in \mathbb{R}^{6}, \quad i=1,\ldots,N$$

    \item \textbf{Exp 1.3 (NoUpright)}:
        Baseline，不使用upright向量
        $$\mathbf{f}_{\text{final}} = \mathbf{f}_{\text{global}} \in \mathbb{R}^{1024}$$
\end{itemize}

\subsection{网络架构}

所有实验均基于PointNet++架构：

\begin{enumerate}
    \item \textbf{Set Abstraction层1}: $10000 \rightarrow 512$ 点，MLP $[64, 64, 128]$
    \item \textbf{Set Abstraction层2}: $512 \rightarrow 128$ 点，MLP $[128, 128, 256]$
    \item \textbf{Set Abstraction层3}: $128 \rightarrow 1$ 点 (global pooling)，MLP $[256, 512, 1024]$
    \item \textbf{分类器MLP}: $1024 \rightarrow 512 \rightarrow 256 \rightarrow 128 \rightarrow 5$
\end{enumerate}

总参数量约为 $1.5$M。

\subsection{训练设置}

\begin{table}[h]
\centering
\caption{训练超参数}
\label{tab:training_params}
\begin{tabular}{ll}
\hline
\textbf{超参数} & \textbf{值} \\
\hline
优化器 & AdamW \\
学习率 & $1 \times 10^{-3}$ \\
权重衰减 & $1 \times 10^{-4}$ \\
学习率调度 & Cosine Annealing \\
Batch大小 & 16 \\
训练轮数 & 200 \\
损失函数 & Cross Entropy (with class weights) \\
数据增强 & 随机旋转 ($\theta \in [0, 2\pi)$ 围绕Y轴) \\
\hline
\end{tabular}
\end{table}

\subsection{实验结果}

\begin{table}[h]
\centering
\caption{消融实验结果对比}
\label{tab:ablation_results}
\begin{tabular}{lccccc}
\hline
\textbf{模型} & \textbf{Upright输入} & \textbf{参数量} & \textbf{Train Acc} & \textbf{Val Acc} & \textbf{Test Acc} \\
\hline
Exp 1.1 & Global Concat & 1.50M & -- & -- & -- \\
Exp 1.2 & Point Concat & 1.50M & -- & -- & -- \\
Exp 1.3 & No Upright & 1.50M & -- & -- & -- \\
\hline
\end{tabular}
\end{table}

\subsubsection{Per-class准确率}

\begin{table}[h]
\centering
\caption{各类别准确率 (测试集)}
\label{tab:per_class_accuracy}
\begin{tabular}{lccccc}
\hline
\textbf{模型} & $K=-1$ & $K=0$ & $K=1$ & $K=2$ & $K=4$ \\
\hline
Exp 1.1 & -- & -- & -- & -- & -- \\
Exp 1.2 & -- & -- & -- & -- & -- \\
Exp 1.3 & -- & -- & -- & -- & -- \\
\hline
\end{tabular}
\end{table}

\subsection{混淆矩阵}

\begin{table}[h]
\centering
\caption{混淆矩阵 (最佳模型，测试集)}
\label{tab:confusion_matrix}
\begin{tabular}{c|ccccc}
\hline
& \multicolumn{5}{c}{\textbf{预测}} \\
\textbf{真实} & $K=-1$ & $K=0$ & $K=1$ & $K=2$ & $K=4$ \\
\hline
$K=-1$ & -- & -- & -- & -- & -- \\
$K=0$ & -- & -- & -- & -- & -- \\
$K=1$ & -- & -- & -- & -- & -- \\
$K=2$ & -- & -- & -- & -- & -- \\
$K=4$ & -- & -- & -- & -- & -- \\
\hline
\end{tabular}
\end{table}

\subsection{讨论}

\subsubsection{关键发现}
\begin{enumerate}
    \item \textbf{Upright向量的重要性}: (待填写实验结果分析)
    \item \textbf{输入方式的影响}: (待填写GlobalConcat vs PointConcat对比)
    \item \textbf{类别不平衡处理}: 使用class weights有效提升了少数类别的准确率
\end{enumerate}

\subsubsection{误分类分析}
(待填写误分类案例分析)

\subsection{消融实验小结}

根据实验结果，我们选择 \textbf{[最佳模型名称]} 作为最终的对称性分类网络，该网络在测试集上达到了 \textbf{XX.X\%} 的准确率。

% 引用示例
% 如需引用表格：如表~\ref{tab:ablation_results}所示...
% 如需引用章节：详见第~\ref{sec:phase2}节...
